\documentclass[11pt,a4paper]{article}
\usepackage[a4paper,margin=22mm]{geometry}
\usepackage{lmodern}
\usepackage{microtype}
\usepackage[hyphens]{url}
\usepackage{hyperref}
\hypersetup{colorlinks=true,urlcolor=blue,pdftitle={Cover Letter: RF3 Input-Channel Audit},pdfauthor={Cheng Ding}}
\setlength{\parindent}{0pt}
\setlength{\parskip}{7pt}
\pagestyle{empty}

\begin{document}

Cheng Ding\\
School of Life Science and Technology\\
ShanghaiTech University, Shanghai, China\\
\href{mailto:dingcheng2024@shanghaitech.edu.cn}{dingcheng2024@shanghaitech.edu.cn}\\
ORCID: \href{https://orcid.org/0009-0001-7221-9066}{0009-0001-7221-9066}

6 November 2026

Dear Editors,

Please consider the manuscript \textbf{``An input-channel audit of RoseTTAFold3: MSA writing, mutation-effect prediction, and conformational-state selection''} for publication as a \textbf{Research Article} in \textit{PLOS Computational Biology}, in the \textbf{Biological Macromolecules} section.

The manuscript addresses a computational structural biology question created by recent AlphaFold3-style architectures: what happens to evolutionary information after the explicit multiple sequence alignment (MSA) representation has been discarded? We audit this input channel in RoseTTAFold3 across three connected levels---pair-representation geometry, deep mutational scanning, and conformational-state selection---using one fixed checkpoint and inference configuration. This multilevel design links an architectural mechanism to both mutation-sensitive functional readouts and generated three-dimensional structures.

We find that evolutionary information written into the pair representation is persistent and functionally useful: a single MSA recycle captures approximately 75\% of the geometric compression and 85\% of the mutation-effect gain obtained with the full inference protocol. Controlled alignment perturbations show that site-wise conservation supplies most of the biological gain, with covariation providing an additional contribution. The mutation-effect advantage persists under train-only standardization, independent test evaluation, contiguous sequence-block cross-validation, and lower-capacity readouts.

We also find that evolutionary alignments and structural templates act as distinct, target-dependent priors. Analysis of two-state proteins and 480 explicitly sampled structures shows that MSA effects on conformational selection are consistent with family consensus, whereas structural templates can suppress mutation sensitivity. Their combined effects do not follow a universal hierarchy: native alignments can preserve or restore a native fold even when opposed by an alternative-fold template. Several template-conditioned distogram preferences disagree with the generated coordinates, demonstrating that conformational audits should include explicit structure generation. Together, these findings provide a new mechanistic account of how evolutionary and structural priors shape both mutation sensitivity and conformational choice in RoseTTAFold3.

The analysis code, evaluation protocols, numerical outputs, and figure-generation scripts are publicly available at \url{https://github.com/Lizon1c/rf3-msa-template-audit}. The sampled coordinates and analysis-ready feature archives are available at \url{https://doi.org/10.5281/zenodo.23859152}. The submission includes two main figures and separate supporting information comprising S1 Text, S1--S2 Figs, and S1--S4 Tables.

This manuscript is original, is not published or under consideration elsewhere, and has been approved by the author. The author declares no competing interests. No external funding was received for this work. Generative AI assistance used for manuscript-text unification and editing is disclosed in the Acknowledgments; the author verified the numerical values, citations, and scientific statements and takes full responsibility for the content.

Thank you for considering this manuscript.

Sincerely,

Cheng Ding

\end{document}
